% !TEX root = ../main.tex

\chapter{Kesimpulan dan Saran}
\label{chap:kesimpulan}


\section{Kesimpulan}
\label{sec:kesimpulan}

Berdasarkan kajian analitik dan penyelesaian numerik yang telah dilakukan,
diperoleh kesimpulan sebagai berikut:

\begin{enumerate}
    \item Persamaan geodesik pada metrik Kerr berhasil diturunkan
          secara eksplisit dari Lagrangian
          $\mathcal{L} = \frac{1}{2}\,g_{\mu\nu}\,\dot{x}^\mu\,\dot{x}^\nu$
          melalui metode Euler-Lagrange.
          Dua konstanta gerak (energi spesifik $E$ dan momentum sudut $L_z$)
          diperoleh dari simetri stasioner dan aksial metrik Kerr.
          Persamaan transportasi paralel tensor spin
          $dS^\mu/d\tau = -\Gamma^\mu_{\nu\rho}\, S^\nu\, u^\rho$
          diformulasikan dalam bentuk sistem matriks $4 \times 4$
          yang dapat diintegrasikan secara numerik bersama dengan
          persamaan geodesik sebagai sistem 12 ODE orde satu.

    \item Metrik Kerr berhasil direduksi ke aproksimasi medan lemah
          melalui ekspansi Taylor terhadap parameter
          $\epsilon_1 = r_s/r \approx 10^{-9}$ dan $\epsilon_2 = a/r \approx 10^{-7}$.
          Komponen campuran $g_{t\phi} = 2GJ\sin^2\!\theta/(c^3 r)$
          diidentifikasi sebagai sumber tunggal efek \textit{frame-dragging}.
          Melalui formulasi gravitomagnetik, ekspresi analitik
          laju presesi Lense-Thirring diperoleh sebagai
          $\Omega_{\text{LT}} = GJ/(2c^2 r^3)$
          untuk orbit polar.
          Substitusi parameter Bumi menghasilkan
          $\Omega_{\text{LT}} \approx 41$~mas/yr
          dan $\Omega_{\text{geodetik}} \approx 6600$~mas/yr,
          yang berada dalam orde yang sama dengan data eksperimen GP-B.

    \item Sistem persamaan diferensial yang dihasilkan
          berhasil diselesaikan secara numerik menggunakan
          metode Runge-Kutta orde 4/5 (RK45).
          Hasil simulasi menunjukkan laju presesi geodesik $\sim 6600$~mas/yr
          dan \textit{frame-dragging} $\sim 37$--$41$~mas/yr,
          konsisten dengan data eksperimen \textit{Gravity Probe B}
          ($6602 \pm 18$~mas/yr dan $37{,}2 \pm 7{,}2$~mas/yr).
          Analisis sensitivitas mengonfirmasi relasi linier
          $\Omega_{\text{LT}} \propto a$
          dan lenyapnya efek pada limit $a = 0$ (Schwarzschild).
          Ketiga besaran diagnostik (normalisasi empat-kecepatan,
          ortogonalitas spin-kecepatan, dan kekekalan norma spin)
          menunjukkan kestabilan numerik yang baik sepanjang integrasi.
\end{enumerate}

Secara keseluruhan, konsistensi antara derivasi analitik,
hasil numerik, dan data eksperimen \textit{Gravity Probe B}
mengonfirmasi validitas penurunan dari metrik Kerr
hingga ekspresi presesi Lense-Thirring
yang dikembangkan dalam penelitian ini.


\section{Saran}
\label{sec:saran}

Berdasarkan hasil penelitian ini, beberapa saran
untuk pengembangan lebih lanjut adalah sebagai berikut:

\begin{enumerate}
    \item Perluasan model orbit dari orbit sirkular ideal
          ke orbit dengan eksentrisitas non-nol
          untuk mengevaluasi pengaruh eksentrisitas
          terhadap laju presesi.

    \item Penggunaan metrik Kerr-Newman
          (massa berotasi dan bermuatan listrik)
          untuk mengkaji efek muatan listrik
          terhadap struktur ruang-waktu dan presesi.

    \item Implementasi integrator simplektik
          (misalnya Störmer-Verlet atau metode dari pustaka FANTASY)
          untuk meningkatkan kestabilan integrasi jangka panjang
          dan melestarikan struktur Hamiltonian sistem.

    \item Validasi terhadap data eksperimen satelit geodesik lain
          seperti LAGEOS, LARES, dan misi masa depan LARES-2
          untuk memperluas rentang verifikasi.

    \item Pengembangan analisis orde tinggi
          (koreksi post-Newtonian orde dua dan efek kuadrupol gravitasi Bumi)
          untuk meningkatkan presisi prediksi teoretik.
\end{enumerate}
